\section{Introduction}

A robot picking up an unfamiliar object, or a legged robot stepping onto ice, must adapt to physics it does not
know from a few seconds of interaction. A popular recipe trains a sequence model across randomised physics and lets it
adapt \emph{in context}, from its own history, with no gradient steps at deployment~\citep{kumar2021rma,duan2016rl2,forgione2023from,busetto2024one}.
These models work, but it is not clear what they compute. Do they identify the plant and then plan? Do they
implement a classical estimator? Do they explore?

For linear-Gaussian dynamics these questions have exact answers, so they can be tested. Transformers trained on
random linear systems are known to \emph{behave} like the Kalman filter or least
squares~\citep{balim2023can,akram2026transformers}. They are also known to be \emph{able to represent} such
estimators~\citep{goel2024can,cole2025incontext,vladarean2026learning}. What has not been measured is what is inside
a trained in-context system-identifier, or what an in-context \emph{controller} does about the dual effect. The second
question matters because, under exploration-free feedback, closed-loop data do not identify the plant~\citep{feldbaum1960dual,astrom1995adaptive,ljung1999system}.
\citet{akram2026transformers} list the first question as future work.

We study both questions in a pre-registered design. Hypotheses, metrics, seeds and tests were committed before the
main experiments, and deviations are logged. The design combines exact optimal references, causal interventions and
three settings of increasing difficulty: linear prediction, MuJoCo contact physics~\citep{todorov2012mujoco}, and adaptive control.
Our contributions:
\begin{itemize}
\item \textbf{Theory (\S\ref{sec:theory}).} An exact Bayes-risk identity for in-context prediction under any
non-anticipating policy, with a parameter-free context-length law $t\,\mathcal{R}^\star_t\to np\sigma^2$. A
closed-loop non-identifiability identity: the information gained along unexcited directions equals the excitation energy, and is
exactly zero under pure feedback. A two-layer construction that gathers the exact sufficient statistics, and a
first-order account of how a near-uniform attention head can apply a parameter estimate that is stored only in its
key--value memory (Remark~\ref{rem:fw}).
\item \textbf{What the trained predictor computes (\S\ref{sec:expa}).} It is Bayes-optimal while data are scarce, and
then it settles at an approximation floor. Our pre-registered mechanism, a linearly readable and causally used parameter estimate in the
residual stream, is \emph{rejected}: the probe reaches $R^2{=}\ARtheta$, and steering and subspace patching have no effect. Exploratory analyses show instead that the estimate lives
in the \emph{key--value memory} of the last two blocks. It is readable from the fast weights ($R^2{=}\AfwR$), it equals the predictor's
query Jacobian ($R^2{=}\AjacBayes$ against the Bayes estimate), and patching only context keys and values transfers
$\AkvThreeFour$ of the behaviour.
\item \textbf{Contact physics (\S\ref{sec:expb}).} On a pushed block with Coulomb friction the transformer beats
physics-feature RLS by $\BKratio\times$, and on a pendulum it beats RLS even though RLS has the true model class. Physical parameters follow the same pattern:
they are weakly readable from activations ($1/m$: $R^2{=}\BprobePinvm$), but strongly readable from the key--value memory ($\BfwPinvm$) and from the model's
input sensitivity ($r{=}\BfuncP$).
\item \textbf{In-context control (\S\ref{sec:expc}).} A transformer controller trained by imitation and then on cost discovers excitation (\Hfived\ nats of information above exploration-free certainty equivalence, of which \CinfoFtBc\ come from cost training), and it matches dithered certainty equivalence and Thompson sampling in median cost. Its failure rate, however, is \CTFfail\% against \CCEDfail\%, concentrated on open-loop-unstable plants, and cost fine-tuning makes those failures more severe. The controller stores its estimate of the actuation matrix in the same key--value memory.
\end{itemize}

\section{Related work}

\paragraph{In-context learning of function classes.} Transformers trained on random regression tasks match least
squares~\citep{garg2022what}. Behaviourally they interpolate between ridge and OLS, and probes recover intermediate
quantities in static regression~\citep{akyurek2022learning}. They can implement GD~\citep{vonoswald2023transformers,ahn2023transformers,zhang2023trained},
higher-order iterations~\citep{fu2023transformers} and algorithm selection~\citep{bai2023transformers}, and they
approximate posterior predictives~\citep{muller2022transformers,raventos2023pretraining}. We move these questions to
dynamical, closed-loop data, where the regressors depend on the past and on the policy. There the dual effect has no analogue in i.i.d.\ regression.

\paragraph{In-context learning of dynamical systems.} Meta-trained transformers predict and filter unseen linear and
nonlinear systems~\citep{forgione2023from,piga2023adaptation,busetto2023incontext,balim2023can,akram2026transformers}
and can serve as contextual controllers for a system class~\citep{busetto2024one}. Constructions and approximation
theory exist for Kalman filtering~\citep{goel2024can}, for least-squares learning of LDS with depth
separation~\citep{cole2025incontext}, and for single linear-attention layers~\citep{vladarean2026learning}. All of these are
either behavioural or representational. None of them intervenes on a trained model.

\paragraph{Adaptation from history and in-context RL.} RMA~\citep{kumar2021rma} regresses a history encoding onto
privileged physical parameters. Probing such latents gives mixed results~\citep{chamachot2026evaluating}, and decoding
them can need a learned generative model~\citep{zhang2026interpreting}. Meta-RL agents can learn to
explore~\citep{duan2016rl2,zintgraf2020varibad}; algorithm distillation imitates learning
histories~\citep{laskin2022incontext}; and in-context RL algorithms are provably representable in tabular
MDPs~\citep{liang2026transformers}. In adaptive LQR, certainty equivalence with dithering is
rate-optimal~\citep{mania2019certainty,simchowitz2020naive}. We use these methods as baselines and ask whether a
transformer controller, trained by imitation and then on cost, discovers the excitation they prescribe.

\paragraph{Interpretability tools and fast weights.} We use linear probes~\citep{alain2016understanding} and activation
patching~\citep{zhang2024towards,heimersheim2024how}. Our probes are always paired with interventions, and the probe
targets are \emph{estimator} quantities, which are distinct from the ground truth. The view of attention as reading and writing
\emph{fast weights}~\citep{ba2016using,schlag2021linear}, and its dual form as implicit gradient
descent~\citep{dai2023why,vonoswald2023transformers}, motivates our key--value analysis. To our knowledge, it has not
been used to localise a physical parameter estimate in a trained model.
